\section{Введение}
\subsection*{Цель} \addcontentsline{toc}{subsection}{Цель}

Цель работы --- изучить динамику систем, совершающих разрывные колебания, на примерах мультивибратора, триггера и кипп-реле: освоить метод разрывных колебаний, определить параметры автоколебаний мультивибратора, а также границы рабочих режимов триггера и кипп-реле.

\subsection*{Задания} \addcontentsline{toc}{subsection}{Задачи}
\begin{enumerate}
	\item С помощью регулятора 1 перевести «Блок релаксационных колебаний» в режим мультивибратора.
	\begin{enumerate}
		\item[(a)] Зарисовать осциллограммы колебаний напряжения и тока, фазовую плоскость мультивибратора с осциллографа.
		\item[(b)] Измерить период и амплитуду автоколебаний мультивибратора.
	\end{enumerate}
	
	\item С помощью регулятора 1 перевести «Блок релаксационных колебаний» в режим триггера, подсоединить генератор импульсов к соответствующему входу. Подать на вход триггера прямоугольные импульсы.
	\begin{enumerate}
		\item[(a)] Измерить длительность импульса на выходе триггера.
		\item[(b)] Выяснить влияние длительности и амплитуды запускающего импульса на переброс триггера.
		\item[(c)] Измерить минимально и максимально допустимые длительности запускающих импульсов. Измерить зависимость амплитуды и периода снимаемого импульса от амплитуды запускающего.
		\item[(d)] Зарисовать осциллограммы напряжения на выходе триггера и входного сигнала в случае деления частоты на триггере.
	\end{enumerate}
	
	\item Подготовить лабораторную установку к изучению колебаний кипп-реле, подсоединив генератор и осциллограф к соответствующим входу и выходу.
	\begin{enumerate}
		\item[(a)] измерить длительность выходного сигнала кипп-реле.
		\item[(b)] Определить минимальную и максимальную длительность запускающего импульса, при которой схема работает как кипп-реле, т.е. длительность выходного сигнала не зависит от длительности входного импульса
		\item[(c)] Выяснить, при каких амплитудах входного импульса спусковая схема запускается и работает как кипп-реле.
	\end{enumerate}
\end{enumerate}

\subsection*{Приборы и оборудование} \addcontentsline{toc}{subsection}{Приборы и оборудование}
\begin{itemize}
	\item Осциллограф GOS-620FG;
	\item Генератор Г5-26;
	\item Блок установки -- самодельный.
\end{itemize}